\documentclass[12pt]{article}
\usepackage[T1]{fontenc}
\usepackage[utf8]{inputenc}
\usepackage{lmodern}
\usepackage{textcomp}
\usepackage{enumitem}
\usepackage{geometry}
\geometry{margin=1in}
\begin{document}
\begin{center}
Answer Key - Physics - Graduate Level Assessment 10 \
\small (Answers are ordered exactly as in the question paper.)
\end{center}

\section*{Multiple Choice}
\begin{enumerate}[label=\arabic*.]
\item \textbf{Answer:} A
\\ \textit{Options:} A) mechanical waves; B) sound waves; C) all waves; D) None of these; E) radio waves
\\ \textit{Note:} Polarization is a property of mechanical waves because these waves can oscillate in specific planes, allowing for the separation of wave vibration directions, which is a characteristic not found in other types of waves such as sound waves.
\item \textbf{Answer:} A
\\ \textit{Options:} A) Sound; B) Ultraviolet waves; C) Infrared waves; D) Microwaves; E) X-rays
\\ \textit{Note:} Sound does not belong in the family of electromagnetic waves because it is a mechanical wave that requires a medium to propagate, whereas electromagnetic waves can travel through a vacuum without a medium.
\item \textbf{Answer:} A
\\ \textit{Options:} A) oval; B) spiral; C) elliptical; D) All of these; E) rectangular
\\ \textit{Note:} The gravitational force on a satellite remains constant in an oval orbit because, despite variations in distance from the central body, the force is determined by the mass of the central body and the satellite, maintaining a consistent interaction throughout the orbit.
\item \textbf{Answer:} A
\\ \textit{Options:} A) [(3 f\_2) / 2]; B) (f\_1 + f\_2) / 2; C) (2 f\_2) / 3; D) 2 f\_2; E) (3 f\_2) / 3
\\ \textit{Note:} The correct answer is derived from the lens combination formula, which states that the effective focal length (F) of two thin lenses separated by a distance (d) is given by 1/F = 1/f\_1 + 1/f\_2 - d/(f\_1*f\_2), and substituting f\_1 = 3 f\_2 and d = 2 f\_2 leads to the result F = (3 f\_2)/2.
\item \textbf{Answer:} C
\\ \textit{Options:} A) Magnetic attraction; B) vacuum polarization; C) interactions with the ionic lattice; D) Gravitational pull; E) Electrostatic attraction
\\ \textit{Note:} The correct answer is based on the fact that, in BCS theory, Cooper pairs form when electrons interact with phonons generated by the vibrations of the ionic lattice, leading to an attractive force that enables superconductivity.
\end{enumerate}

\section*{True / False}
\begin{enumerate}[label=\arabic*.]
\setcounter{enumi}{5}
\item \textbf{Answer:} False
\\ \textit{Options:} A) True; B) False
\\ \textit{Note:} The statement is false because relativistic equations for time dilation and length contraction apply at speeds approaching the speed of light, but not at speeds greater than the speed of light, which is prohibited by the theory of relativity.
\item \textbf{Answer:} False
\\ \textit{Options:} A) True; B) False
\\ \textit{Note:} The statement is false because, according to Ohm's law (V = IR), the current through the light bulb is actually 0.5 amperes, calculated by dividing 120 volts by 240 ohms.
\item \textbf{Answer:} False
\\ \textit{Options:} A) True; B) False
\\ \textit{Note:} The total spin quantum number of the electrons in the ground state of neutral nitrogen is 1/2, as nitrogen has three unpaired electrons with a total spin of 3/2, not the stated value of 3.
\item \textbf{Answer:} False
\\ \textit{Options:} A) True; B) False
\\ \textit{Note:} Jupiter's internal heat is primarily generated by gravitational contraction and the slow release of heat from its formation, rather than nuclear fusion, which is a process that occurs in the cores of stars.
\item \textbf{Answer:} False
\\ \textit{Options:} A) True; B) False
\\ \textit{Note:} The buoyant force on the balloon remains constant as it sinks deeper into the lake because it depends only on the volume of fluid displaced, which does not change with depth, despite the increase in water pressure.
\end{enumerate}

\section*{Long Form Response}
\begin{enumerate}[label=\arabic*.]
\setcounter{enumi}{10}
\item \textbf{Answer:} In designing a two-mile-long steel bridge, it is essential to account for thermal expansion due to temperature variations ranging from -40^\ensuremath{\circF} to +110^\ensuremath{\circF}. The total temperature change is 150^\ensuremath{\circF}, and steel expands approximately 6.5 x 10^-6 per degree Fahrenheit. Applying the formula for linear expansion, the total expansion can be calculated as: Total Expansion = Original Length x Coefficient of Expansion x Temperature Change, which results in: 2 miles x 6.5 x 10^-6/^\ensuremath{\circF} x 150^\ensuremath{\circF} = 0.003 miles. This expansion of 0.003 miles must be considered in the structural design to prevent stress and potential failure, necessitating the inclusion of expansion joints and flexible supports to accommodate such movement.
\\ \textit{Note:} The answer correctly applies the linear expansion formula to calculate the total thermal expansion of a two-mile-long steel bridge due to a temperature variation of 150°F, demonstrating an understanding of the material properties of steel and the practical implications for bridge design.
\item \textbf{Answer:} The total charge for a mass of 75.0 kg of electrons can be calculated using the formula \( Q = n \cdot e \), where \( n \) is the number of electrons and \( e \) is the elementary charge (approximately \( 1.6 \times 10^{-19} \) C). First, we determine the number of electrons \( n \) by using the mass of an electron (approximately \( 9.11 \times 10^{-31} \) kg): \( n = \frac{75.0 \, \text{kg}}{9.11 \times 10^{-31} \, \text{kg/electron}} \approx 8.23 \times 10^{31} \, \text{electrons} \). Consequently, the total charge is \( Q = (8.23 \times 10^{31}) \cdot (-1.6 \times 10^{-19} \, \text{C}) \approx -1.32 \times 10^{13} \, \text{C} \). This substantial negative charge has significant implications in electrostatics, such as influencing electric fields and forces in charged systems, and it also raises considerations in real-world applications, including the design of electronic devices and understanding charge transport in materials.
\\ \textit{Note:} correct because it accurately applies the relationship between mass and number of electrons to find the total charge, demonstrating the fundamental principles of electrostatics and the significant implications of such a large negative charge in real-world contexts like electrical engineering and materials science.
\end{enumerate}

\vfill
\noindent\textit{End of Answer Key}
\end{document}